\documentclass[11pt]{article}
\usepackage[margin=1in]{geometry}
\usepackage[T1]{fontenc}
\usepackage{lmodern}
\usepackage{amsmath,amssymb,amsthm}
\usepackage{booktabs}
\usepackage{microtype}
\usepackage{multicol}
\usepackage[inline]{enumitem}
\setlist{nosep}
\setlength{\abovedisplayskip}{4pt}
\setlength{\belowdisplayskip}{4pt}
\setlength{\abovedisplayshortskip}{2pt}
\setlength{\belowdisplayshortskip}{2pt}
\usepackage[hidelinks]{hyperref}

\newtheorem{theorem}{Theorem}
\newtheorem{lemma}[theorem]{Lemma}
\newtheorem{proposition}[theorem]{Proposition}
\newtheorem{corollary}[theorem]{Corollary}
\theoremstyle{definition}
\newtheorem{definition}[theorem]{Definition}
\theoremstyle{remark}
\newtheorem{observation}[theorem]{Observation}
\newtheorem{remark}[theorem]{Remark}

\newcommand{\PH}{P^{\mathrm H}}
\newcommand{\PV}{P^{\mathrm V}}
\newcommand{\Zs}{\mathbb{Z}_7^5}
\newcommand{\boxt}{\boxtimes}

\title{The shape of a private pair redistributes its cost:\\ private pairs as a property of the code}
\author{Artem Oktiabrev\thanks{Independent researcher, Ukraine.
\href{https://orcid.org/0009-0003-3626-2002}{ORCID 0009-0003-3626-2002},
\texttt{aoktyabrev@gmail.com}.}}
\date{8 October 2026}

\begin{document}
\maketitle

\begin{abstract}
% The abstract of the second note.  Every number is checked by scripts/w2_checknums.py.
The recent lower bounds on the Shannon capacity $\Theta(C_7)$ all rest on gadgets in
$C_7^{\boxt 5}$: a code, a set of \emph{private pairs} (a vertex $q$ outside the code with a
single neighbour $r$ in it), and an auxiliary set that must avoid a forbidden region and is
rewarded for avoiding a penalised one. A companion note proves that the Polak--Schrijver code
admits exactly eight private pairs. This note explains why the number of pairs is not, by
itself, a lever. The two regions a pair creates are the intersection and the union of two
closed neighbourhoods, so their sizes always add up to $2\cdot 3^d$: the shape of a pair moves
cost between the forbidden and the penalised region but cannot lower it. We list the five shapes
in $C_7^{\boxt 5}$ with their costs and the largest independent set we could build carrying each;
show that every $367$-word code we can reach carries only the two-coordinate shape; that up to
automorphism every nine-pair code we can reach has the same sixteen forbidden regions; and that
against all of them the best auxiliary set has $366$ words, one short, with a status we state
precisely: it is $3$-opt optimal, and not proved maximum.

\end{abstract}

% The body of the second note.  Every number in this file is checked by
% scripts/w2_checknums.py against results/json.
\section{Context}
\label{sec:intro}

Since July 2026 every improvement of the lower bound on $\Theta(C_7)$
\cite{itty2026,gao2026,bpz2026,tandon2026} starts from one five-dimensional gadget
\cite{gao2026}: the $367$-word Polak--Schrijver code $I_0$ \cite{polak2019}, eight private
pairs, two transversals of their endpoints, and an auxiliary set $X$. The recursions built on
it are most sensitive to the number $t$ of private pairs, and a companion note
\cite{oktiabrev2026a} proves that for $I_0$ this number cannot be raised: $t^*(I_0)=8$. That
leaves the obvious next question --- another code with more pairs --- and this note answers why
that route does not pay as directly as the sensitivity suggests. Together with
\cite{oktiabrev2026a} it accounts for the route ``more private pairs'' on every code we hold:
on $I_0$ there is no ninth pair, and on the codes that have nine or ten the best auxiliary set
falls one word short at nine pairs and two at ten, against forbidden regions that are fixed by
the geometry of the pairs rather than by the code. That is a statement about
the codes we hold, not about all codes; Section~\ref{sec:open} says what would change it.

Only one statement below is proved by hand, and it is elementary. Everything else is a finite
computation, exhaustive where we say so and a search where we say so; the status is given with
each item, as in \cite{oktiabrev2026a}.

\section{Definitions}
\label{sec:defs}

Vertices of $C_7^{\boxt d}$ are the elements of $\mathbb{Z}_7^d$; two distinct vertices are
confusable iff their circular distance is at most $1$ in every coordinate. $N(S)$ is the
\emph{closed} neighbourhood of a set $S$, as in \cite{gao2026,tandon2026}; for a single
vertex $|N(v)|=3^d$. A pair $(r,q)$ with $r\in I$, $q\notin I$ is a \emph{private pair} for an
independent set $I$ if $N(q)\cap I=\{r\}$ \cite{gao2026,tandon2026}. A gadget
\cite{gao2026} has a code $I$, private pairs $(r_i,q_i)$, $1\le i\le t$, with disjoint
endpoint sets, two independent complementary transversals $\PH,\PV$ (one endpoint of every
pair in each), and an independent $X$ with $X\cap N(\PH)\cap N(\PV)=\emptyset$; its profile
counts $o=|X\setminus(N(\PH)\cup N(\PV))|$. We name the two sets in that definition:
\[
  F=N(\PH)\cap N(\PV)\quad\text{(forbidden)},\qquad
  U=N(\PH)\cup N(\PV)\quad\text{(penalised)}.
\]
$X$ must avoid $F$, and every word of $X$ in $U$ is lost from $o$. The \emph{shape} of a pair is
$\delta=q-r\in\{0,\pm1\}^d\setminus\{0\}$; $k$ is its number of non-zero coordinates.

\section{Cost is conserved}
\label{sec:identity}

\begin{lemma}
\label{lem:conserve}
Let $n\ge3$ and let $\PH,\PV$ be any two sets of vertices of $C_n^{\boxt d}$. Then
\[
  |F|+|U| \;=\; |N(\PH)|+|N(\PV)| \;\le\; 3^d\,(|\PH|+|\PV|),
\]
with equality iff the closed neighbourhoods of the vertices inside each of $\PH$ and $\PV$ are
pairwise disjoint. In particular, for a single pair, $\PH=\{r\}$ and $\PV=\{q\}$,
\[
  |N(r)\cap N(q)|+|N(r)\cup N(q)| \;=\; 2\cdot 3^d .
\]
\end{lemma}

\begin{proof}
$|A\cap B|+|A\cup B|=|A|+|B|$ for any finite sets, with $A=N(\PH)$, $B=N(\PV)$; and
$|N(S)|\le\sum_{s\in S}|N(s)|=3^d|S|$, since for $n\ge3$ the three values $a-1,a,a+1$ of
every coordinate are distinct.
\end{proof}

The lemma is inclusion--exclusion, and it holds for any two vertices, private or not; it is
stated because of what it implies for gadgets. A pair contributes $N(r)\cap N(q)$ to $F$ and
$N(r)\cup N(q)$ to $U$, and the two always add up to $2\cdot 3^d$, which is $486$ in
$C_7^{\boxt 5}$. \emph{The shape of a pair redistributes its cost between the forbidden and the
penalised region; it does not lower it.} A shape that makes $F$ small makes $U$ large by the
same amount, and $U$ is what $o$ is paid from.

\begin{lemma}
\label{lem:shape}
For $n\ge4$ and a pair of shape $k$, $|N(r)\cap N(q)|=3^{d-k}\,2^{k}$.
\end{lemma}

\begin{proof}
Coordinatewise, $\{a-1,a,a+1\}\cap\{a+\delta-1,a+\delta,a+\delta+1\}$ has $3$ elements if
$\delta=0$ and $2$ if $\delta=\pm1$, for $n\ge4$; closed neighbourhoods in the strong product
are products of the coordinate neighbourhoods.
\end{proof}

Both lemmas were also checked by direct count for every shape with $d=3,4,5$ and
$n=3,4,5,7,9$ ($n=3$ satisfies Lemma~\ref{lem:conserve} and, as it must, not
Lemma~\ref{lem:shape}). For families the inequality is strict in practice: the nine-pair
configurations of Section~\ref{sec:regions} have $|F|+|U|$ between $4050$ and $4080$ against
$2\cdot9\cdot3^5=4374$.

\section{The five shapes}
\label{sec:shapes}

The $242$ shapes in $C_7^{\boxt 5}$ fall into five orbits under the stabiliser of a vertex in
$(D_7)^5\rtimes S_5$, one per $k$. The last column is a finite computation, exhaustive over
its inputs: for every word $r$ of every code we hold and of the swap plateau of
Section~\ref{sec:regions}, and every shape, put $q=r+\delta$, delete the other code words
confusable with $q$, and add back an exact maximum independent set of the freed vertices that
keep $q$ private. Every set reported is verified, and its pair is re-checked to be private by
a second program.

\begin{center}\small
\begin{tabular}{@{}crrrcr@{}}
\toprule
$k$ & shapes & $|N(r)\cap N(q)|$ & $|N(r)\cup N(q)|$ & covered cells in box$(q)$ & largest set found \\
 & & (forbidden) & (penalised) & & \\
\midrule
1 & 10 & 162 & 324 & 16 & 366 \\
2 & 40 & 108 & 378 & 8  & 367 \\
3 & 80 & 72  & 414 & 4  & 366 \\
4 & 80 & 48  & 438 & 2  & 366 \\
5 & 32 & 32  & 454 & 1  & 366 \\
\bottomrule
\end{tabular}
\end{center}

The cost in $F$ \emph{falls} with $k$; the two-coordinate shape is in the middle, not at the
minimum. Every shape is realisable as a private pair, the expensive $k=1$ included --- all seven
candidates of the Mathew--\"Osterg{\aa}rd $350$-word set \cite{mathew2017} have $k=1$ --- but
in the sets we can build, shapes other than $k=2$ occur only one word below $367$.

The column of covered cells explains why. A word $w$ occupies the box $w+\{0,1\}^5$ of $32$
cells, and two vertices are non-confusable exactly when their boxes are disjoint; a private
$q$ meets the code's boxes only inside box$(r)$, in $2^{5-k}$ cells. A pair of shape $k\ge3$
therefore needs a box with at most $4$ covered cells. In every one of the $367$-word codes we
can reach, every box outside the code has at least $5$.

\section{The count is a property of the code}
\label{sec:count}

Counting the candidates --- the vertices with exactly one neighbour in the code --- for every
code we hold gives a statistic that, as far as we can find, no paper reports. The counts do not
follow size: the linear $343$-word code \cite{baumert1971} has none, the $350$-word set of
\cite{mathew2017} has seven, the printed $367$-word code has eight, and an independent rebuild
of it that differs in two words has six. Among the $367$-word codes the counts we see run from
$5$ to $10$, every candidate has shape $k=2$, and for a fixed code the forbidden region varies
over the admissible colourings by at most $3.21\%$.

\section{Sixteen forbidden regions}
\label{sec:regions}

A private pair is also a move: $I\mapsto I-r+q$ keeps size and independence, and it is the
exchange used in \cite{tandon2026}. Closing every $367$-word code we hold under these moves,
and under exhaustive exchanges of two and three words, gives a finite plateau that falls into
$51$ classes under $\mathrm{Aut}(C_7^{\boxt 5})$. Each connected component is a product of
five independent three-state switches, $3^5=243$ codes, and the five switches sit on the five
cyclically adjacent pairs of coordinates, permuted by an automorphism of order $5$ that cycles
the coordinates. (The plateau first appeared with $486$ codes, two components; that $486$ also
equals $2\cdot3^5$ is a coincidence: its $3$ counts states of a switch and its $2$ counts the
components our starting codes happened to lie in --- a third start adds a third.)

Two classes have nine candidates: the published auxiliary set of \cite{gao2026} read as a code,
and one more; one class has ten. An automorphism $g$ moves the forbidden region to $g\cdot F$
and leaves $\alpha(C_7^{\boxt5}\setminus F)$ unchanged, so the variables are the code up to
$\mathrm{Aut}$ and the colouring. All $208$ configurations with $t\ge9$ give, up to
automorphism, exactly $16$ forbidden regions at $t=9$ and $4$ at $t=10$, with
\[
  1098\le|F|\le1128\ (t=9),\qquad 1240\le|F|\le1268\ (t=10),
\]
and \emph{each of the sixteen is produced by every nine-pair class}. The forbidden region is a
function of the geometry of the nine pairs, which are the same switches in every code, and not
of the rest of the code. In the nine-pair regions $972=9\cdot108$ vertices come from the pairs
themselves and $126$ to $156$ from cross terms between pairs.

\section{One word short: what is known, and with what status}
\label{sec:gap}

Against each of the $20$ regions, every image under $(D_7)^5\rtimes S_5$ of one representative
of each of the $51$ classes was scored exhaustively: $64{,}538{,}880$ group elements per pair,
$65{,}829{,}657{,}600$ images in all. No image avoids any region; the fewest words of an image
in $F$ is $1$. Every image within three words of admissible was repaired exactly, and the best
results are
\[
  366 \text{ words at } t=9,\qquad 365 \text{ at } t=10 .
\]
All $108$ repaired $366$-word sets are exactly $3$-opt optimal inside $G_F=C_7^{\boxt5}\setminus
F$: no window of at most three words can be traded for one more, decided by counting blockers
rather than by search. Only $2$ or $3$ allowed vertices are blocked by a single word of such a
set.

\emph{Status.} This is not a proof that $367$ words do not fit outside $F$. The cheap upper
bounds cannot see $F$: the Lov\'asz bound gives $401$ whether $F$ is empty or removes $1268$
vertices, because it bounds the whole graph, and a box-volume count is weaker still. An honest
bound needs $\vartheta$ of the induced subgraph on $15{,}679$ to $15{,}709$ vertices, and $F$
breaks the symmetry that would reduce it.

\emph{Why a cheaper shape does not obviously help.} With $|U|$ taken without overlaps, the
fraction of its random share that $X$ may touch at $t=9$ and $367$ words is $0.879$, $0.754$,
$0.688$, $0.651$, $0.628$ for $k=1,\dots,5$; the best auxiliary sets on record reach $0.7678$.
The shapes that would shrink $F$ are the ones that ask most of $o$ (Lemma~\ref{lem:conserve}),
and at $366$ words no shape allows more than $0.68$ of the random share at $t=9$, $13$ or $16$. These are distances to
the best known search, not impossibility results.

\section{What remains open}
\label{sec:open}

\begin{enumerate}
\item \emph{A $367$-word code with a deeper hole}: a box covered in at most $4$ cells, which a
pair of shape $k\ge3$ needs. None of the codes we can reach has one; finding one requires
$367$-word codes from outside the known family, which is the open problem around
$\alpha(C_7^{\boxt 5})$ itself.
\item \emph{An upper bound that sees $F$}, of $\vartheta$ strength on the induced subgraph, which
would turn the one-word gap of Section~\ref{sec:gap} into a theorem or refute it.
\end{enumerate}

\section{Reproducibility}
\label{sec:repro}

All computations are exact integer computations, published with the scripts that produce them,
the predictions registered before each run, the calibrations each estimator had to pass (among
them: the new group sweep reproduces the histogram of \cite{oktiabrev2026a} exactly, the repair
reaches $367$ at $t=8$, the window search recovers $367$ from deliberately damaged sets, and the
automorphism checker rejects a non-automorphism), and a machine reconciliation of every number
in this note against the stored results. They were produced through Claude Code under the
author's direction.

\paragraph{Data and code availability.}
The sets, the scripts that produce every number in this note, the machine reconciliation of
those numbers against the stored results, and the literature check are archived at
\href{https://doi.org/10.5281/zenodo.23234251}{doi:10.5281/zenodo.23234251}~\cite{shannonarchive2}.

\paragraph{Declaration of competing interest.}
The author declares that he has no known competing financial interests or personal relationships
that could have appeared to influence the work reported in this paper.

\paragraph{Funding.}
This research did not receive any specific grant from funding agencies in the public, commercial,
or not-for-profit sectors.

\paragraph{Declaration of generative AI and AI-assisted technologies in the manuscript
preparation process.}
During the preparation of this work the author used Claude and Claude Code (Anthropic) in order
to implement the computations and draft the text. After using this tool/service the author
reviewed and edited the content as needed and takes full responsibility for the content of the
published article.


\footnotesize
\begin{multicols}{2}
\begin{thebibliography}{99}
\setlength{\itemsep}{-1pt}
\setlength{\parsep}{0pt}
\bibitem{oktiabrev2026a} A.~Oktiabrev, The Polak--Schrijver code has exactly eight private
pairs, preprint, Zenodo, 2026.
\href{https://doi.org/10.5281/zenodo.22979509}{doi:10.5281/zenodo.22979509}.
\bibitem{shannonarchive2} A.~Oktiabrev, Sets, verifier and scripts for ``The shape of a private
pair redistributes its cost'' [dataset], Zenodo, v1.2.0, 2026.
\href{https://doi.org/10.5281/zenodo.23234251}{doi:10.5281/zenodo.23234251}.
\bibitem{baumert1971} L.~Baumert, R.~McEliece, E.~Rodemich, H.~Rumsey, R.~Stanley, H.~Taylor,
A combinatorial packing problem, \textit{Computers in Algebra and Number Theory}, American
Mathematical Society, 1971.
\bibitem{bpz2026} P.~Buys, S.~Polak, J.~Zuiddam, Lean-verified lower bounds for the Shannon
capacity of odd cycles, arXiv preprint, 2026. \href{https://doi.org/10.48550/arXiv.2607.29681}{doi:10.48550/arXiv.2607.29681}.
\bibitem{gao2026} Y.~Gao, A recursive construction improving the lower bound on the Shannon
capacity of $C_7$, arXiv preprint, 2026. \href{https://doi.org/10.48550/arXiv.2607.27869}{doi:10.48550/arXiv.2607.27869}.
\bibitem{itty2026} N.~Itty, C.~D.~Rosin, C.~Carstensen, D.~Reichman, Improved lower bounds
for the Shannon capacity of odd cycles, arXiv preprint, 2026. \href{https://doi.org/10.48550/arXiv.2607.21517}{doi:10.48550/arXiv.2607.21517}.
\bibitem{lovasz1979} L.~Lov\'asz, On the Shannon capacity of a graph, \textit{IEEE
Transactions on Information Theory} 25 (1979) 1--7.
\bibitem{mathew2017} K.~A.~Mathew, P.~R.~J.~\"Osterg{\aa}rd, New lower bounds for the Shannon
capacity of odd cycles, \textit{Designs, Codes and Cryptography} 84 (2017) 13--22.
\bibitem{polak2019} S.~C.~Polak, A.~Schrijver, New lower bound on the Shannon capacity of
$C_7$ from circular graphs, \textit{Information Processing Letters} 143 (2019) 37--40.
\bibitem{tandon2026} R.~Tandon, Strengthening recursive constructions for zero-error Shannon
capacity, arXiv preprint, 2026. \href{https://doi.org/10.48550/arXiv.2608.30273}{doi:10.48550/arXiv.2608.30273}.
\end{thebibliography}
\end{multicols}

\end{document}
